\section{\texttt{zoo-bot}}
\label{sec:zoo-bot}

\texttt{zoo-bot} is \Zoo{}'s pure-Go re-write of the \Hanzo{} OpenClaw bot
(at \texttt{\textasciitilde/work/hanzo/bot}). The original is a TypeScript
agentic bot framework that powers content automation, on-chain triggers,
and Discord/Slack integrations across the Hanzo ecosystem. \Zoo{}'s
re-write keeps the surface --- skills, tools, MCP integration, web
hooks --- but re-implements the runtime in Go for size, speed, and
attack surface.

\subsection{Why a re-write}

The OpenClaw original is good. It is also a Node runtime carrying its
TypeScript ecosystem of dependencies: hundreds of npm packages, an
asynchronous event loop, a JIT, and a garbage collector tuned for
front-end-style workloads. For agentic workflows that run in
continuous-loop mode on a server, this is overhead. \texttt{zoo-bot}
trades the front-end ergonomic toolchain for a single static binary.

\subsection{Goals}

\begin{enumerate}[leftmargin=1.4em]
  \item \textbf{Smaller binary.} Single static Go binary. Target $<$50 MB
        stripped, versus the original's roughly 200 MB Node
        runtime + dependencies.
  \item \textbf{Lower memory.} Idle footprint under 80 MB resident,
        versus the original's typical 250+ MB.
  \item \textbf{Faster cold start.} Sub-100 ms startup to first
        accepting connection, versus several seconds for the original.
  \item \textbf{Smaller attack surface.} No npm dependency tree. No
        post-install scripts. No
        transitive supply-chain exposure beyond Go's standard library
        plus a vetted set of direct module dependencies.
  \item \textbf{Drop-in surface.} Agents written for OpenClaw run on
        \texttt{zoo-bot} via a compatibility shim that translates the
        TS skill manifest to Go's plugin descriptor.
\end{enumerate}

\subsection{Architecture}

\texttt{zoo-bot} is organised into three layers:

\begin{itemize}[leftmargin=1.2em]
  \item \textbf{Core.} Event loop, scheduler, MCP server/client,
        tool dispatcher, persistence (SQLite or Postgres).
  \item \textbf{Skills.} Pluggable Go modules implementing the
        \texttt{Skill} interface. Each skill is a small struct with
        \texttt{Plan}, \texttt{Execute}, and \texttt{Observe}
        methods.
  \item \textbf{Adapters.} Connectors to Discord, Slack, Matrix,
        webhook endpoints, and on-chain trigger streams from \Zoo{}
        L1 / \Lux{} chains.
\end{itemize}

\subsection{Use cases}

\begin{itemize}[leftmargin=1.2em]
  \item \textbf{Agentic workflows.} An agent that monitors a per-\LLM{}
        chain attribution graph and posts contributor leaderboards to
        Discord.
  \item \textbf{Content automation.} An agent that watches an A-Chain
        attestation feed and publishes provenance-verified inference
        outputs to a public mirror.
  \item \textbf{On-chain triggers.} An agent that fires when a wildlife
        \NFT{} crosses a price threshold or when a per-\LLM{} chain
        emits a new training-data attestation.
  \item \textbf{Validator-side automation.} An agent that watches the
        \Zoo{} validator's mempool, surfaces anomalies, and pages the
        operator.
\end{itemize}

\subsection{Security model}

Each skill runs in a Goroutine sandboxed by a capability set: filesystem
paths it can read/write, network endpoints it can dial, RPC endpoints
it can call. The sandbox is enforced by a syscall-level hook (Linux
seccomp profile, macOS sandbox profile) that the bot loads at startup.
Skills cannot escape their declared capabilities without a code
modification and a re-deploy.

\subsection{Relationship to OpenClaw}

OpenClaw remains the upstream reference. \texttt{zoo-bot} is a
re-implementation that interoperates with the OpenClaw skill manifest
format, so a skill author can target both. \Hanzo{} continues to
maintain OpenClaw; \Zoo{} maintains \texttt{zoo-bot}. The two share
the MCP wire protocol so a \texttt{zoo-bot} agent can call OpenClaw
tools and vice versa.
